\NativeSlot{53}{36}{%
% p055_s01

\ExamQuestion{2（1）}{7}%
求下列各数的近似数。精确到百分位：\(2.9633\approx(\quad)\)，\allowbreak{}\(249.3370\approx(\quad)\)。\par
}
\NativeSlot{101}{39}{%
% p055_s02

\ExamQuestion{2（2）}{7}%
先改写成用“万”或“亿”作单位的数，再保留一位小数。\(327900=(\quad)\text{万}\approx(\quad)\text{万}\)；\allowbreak{}\(25499270000=(\quad)\text{亿}\approx(\quad)\text{亿}\)。\par
}
\NativeSlot{176}{47}{%
\ExamQuestion{3.}{5}%
\(\square\)里最大能填几？\par\smallskip
\noindent%
% p055_s03
\makebox[.332\linewidth][l]{\(3.\square>3.68\)}%
% p055_s04
\makebox[.332\linewidth][l]{\(2.07>2.\square6\)}%
% p055_s05
\makebox[.332\linewidth][l]{\(3240\text{克}>\square.15\text{千克}\)}\par
\noindent%
% p055_s06
\makebox[.332\linewidth][l]{\(6.25\text{吨}>6\square25\text{千克}\)}%
% p055_s07
\makebox[.332\linewidth][l]{\(0.\square5\text{米}>15\text{厘米}\)}%
% p055_s08
\makebox[.332\linewidth][l]{\(21.5>21.\square9\)}\par
}
\NativeSlot{261}{64}{%
\ExamGroup{四、操作。{\normalfont\kaishu\fontsize{10.8bp}{14bp}\selectfont （共12分）}}

% p055_s09

\ExamQuestion{1.}{5}%
涂色表示下面的小数：\(1.4\)，\allowbreak{}\(1.07\)。每个完整正方形表示1。\par
\BookFigure{figures/p055_s09-01.tex}
}
\NativeSlot{378}{110}{%
% p055_s10

\ExamQuestion{2（1）}{5}%
如下表是4名同学的身高信息。根据线段图回答问题。（1）箭头所表示的是\mbox{（\hspace{1.8em}）}的身高。\par
\par\smallskip\Needspace{4\baselineskip}\begingroup\small\setlength{\tabcolsep}{3pt}
\noindent\begin{tabularx}{\linewidth}{|*{5}{>{\centering\arraybackslash}X|}}\hline
学生 & 乐乐 & 王华 & 可可 & 冬冬\\\hline
身高/米 & 1.32 & 1.19 & 1.45 & \(1.\square4\)\\\hline
\end{tabularx}\par
\endgroup\smallskip

% p055_s11

\ExamQuestion{2（2）}{5}%
如下表是4名同学的身高信息。根据线段图回答问题。（2）在图中标出乐乐和王华的身高。\par

% p055_s12

\ExamQuestion{2（3）}{5}%
如下表是4名同学的身高信息。根据线段图回答问题。（3）冬冬的身高从高到低排第三，他的身高是\mbox{（\hspace{1.8em}）}米。\par
\BookFigure{figures/p055_s10-01.tex}
}
\NativeSlot{566}{63}{%
\ExamGroup{五、解决问题。{\normalfont\kaishu\fontsize{10.8bp}{14bp}\selectfont （共31分）}}

% p055_s13

\ExamQuestion{1.}{7}%
一块正方形菜地的边长缩小到原来的\(\frac1{100}\)之后是\(0.03\)米，原来这块菜地的面积是多少平方米？\par
}
